\documentclass[11pt,a4paper]{article}

\usepackage[margin=2.3cm]{geometry}
\usepackage[T1]{fontenc}
\usepackage{lmodern}
\usepackage{microtype}
\usepackage{xcolor}
\usepackage{booktabs}
\usepackage{longtable}
\usepackage{array}
\usepackage{enumitem}
\usepackage{fancyhdr}
\usepackage{titlesec}
\usepackage{listings}
\usepackage[hidelinks]{hyperref}

\definecolor{brand}{HTML}{3B5BDB}
\definecolor{codebg}{HTML}{F4F5F7}

\titleformat{\section}{\Large\bfseries\color{brand}}{\thesection}{0.6em}{}
\titleformat{\subsection}{\large\bfseries}{\thesubsection}{0.6em}{}

\pagestyle{fancy}
\fancyhf{}
\fancyhead[L]{AgentForge -- User Guide \& User Journey}
\fancyhead[R]{\thepage}
\renewcommand{\headrulewidth}{0.3pt}

\lstset{
  basicstyle=\ttfamily\small,
  backgroundcolor=\color{codebg},
  frame=single,
  rulecolor=\color{codebg!70!black},
  breaklines=true,
  columns=fullflexible,
  keepspaces=true
}

\setlength{\parindent}{0pt}
\setlength{\parskip}{0.55em}
\renewcommand{\arraystretch}{1.25}

\newcolumntype{L}[1]{>{\raggedright\arraybackslash}p{#1}}

\begin{document}

% ------------------------------------------------------------------ title
\begin{titlepage}
  \centering
  \vspace*{3cm}
  {\Huge\bfseries\color{brand} AgentForge\par}
  \vspace{0.6cm}
  {\Large User Guide \& User Journey\par}
  \vspace{1.5cm}
  {\large A no-code platform for building AI-powered workflows\par}
  \vfill
  {\large Document for technical review\par}
  \vspace{0.3cm}
  {\large \today}
  \vspace{2cm}
\end{titlepage}

\tableofcontents
\newpage

% ------------------------------------------------------------------ 1
\section{Overview}

AgentForge is a platform on which users build \textbf{workflows}: blocks joined
together that the computer runs step by step. A workflow can contain AI blocks,
lookups in the company database, and calls to tools such as Slack, GitHub or
Google Sheets. People who do not want to use the editor can run a finished
workflow from a simple page: fill in a few boxes, press one button, read the
result.

\subsection{Key concepts}

\begin{longtable}{@{}L{3.4cm}L{12cm}@{}}
\toprule
\textbf{Term} & \textbf{Meaning} \\
\midrule
\endhead
Workflow & A set of connected blocks, like a recipe the computer follows. \\
Block (node) & One step in a workflow. \\
Run & One execution of a workflow, with its own result and history. \\
Workspace & A team's space. Members have the role owner, editor or viewer. \\
Credential & A saved, encrypted login or API key for a tool. \\
Company data & The company's own database, connected \textbf{read-only}. \\
Copilot & An AI helper that builds, edits, explains and fixes workflows from plain-English requests. \\
\bottomrule
\end{longtable}

\subsection{System components}

\begin{itemize}[nosep]
  \item \textbf{Web app} (Next.js, port 3000): editor, Use page, settings.
  \item \textbf{API} (FastAPI, port 8000): graphs, credentials, files, data sources, audit log, Copilot.
  \item \textbf{Worker} (Python): executes runs through Temporal.
  \item \textbf{Temporal} (port 7233, UI on 8233): durable workflow engine.
  \item \textbf{Database}: SQLite by default, PostgreSQL in Docker. \textbf{Redis} for caching and limits.
\end{itemize}

% ------------------------------------------------------------------ 2
\section{Installing and running locally}

\subsection{Option A: Docker (recommended)}

Prerequisites: Git and Docker Desktop (WSL2 on Windows).

\begin{lstlisting}[language=bash]
git clone https://github.com/jeorgeiiii/Agent_Forge-.git
cd Agent_Forge-
cp .env.example .env
\end{lstlisting}

Edit \texttt{.env} and set:

\begin{itemize}[nosep]
  \item \texttt{INTERNAL\_API\_TOKEN}: any long random string.
  \item \texttt{CREDENTIALS\_MASTER\_KEY}: generate with
  \texttt{python -c "import os,base64;print(base64.urlsafe\_b64encode(os.urandom(32)).decode())"}
  \item \texttt{GROQ\_API\_KEY} (or \texttt{COPILOT\_API\_URL} and \texttt{COPILOT\_API\_KEY}): only needed for the Copilot and AI blocks. Each person should create their own key; do not share keys.
\end{itemize}

\begin{lstlisting}[language=bash]
docker compose -f infra/docker-compose.yml --env-file .env up --build
\end{lstlisting}

\begin{longtable}{@{}L{4cm}L{8cm}@{}}
\toprule
\textbf{Service} & \textbf{URL} \\
\midrule
Web app & \url{http://localhost:3000} \\
API documentation & \url{http://localhost:8000/docs} \\
Temporal UI & \url{http://localhost:8233} \\
\bottomrule
\end{longtable}

\subsection{Option B: Native run (for debugging)}

Prerequisites: Python 3.12+, \texttt{uv}, Node 20+, Docker (for Temporal and Redis).
Use four terminals from the repository root:

\begin{enumerate}[nosep]
  \item Backing services: \texttt{docker compose -f infra/docker-compose.yml up postgres temporal redis} (Postgres is optional because the default \texttt{.env} uses SQLite).
  \item API: \texttt{cd apps/api \&\& uv sync \&\& uv run uvicorn app.main:app --reload --port 8000}
  \item Worker: \texttt{cd apps/worker \&\& uv sync \&\& uv run python main.py}
  \item Web: \texttt{cd apps/web \&\& npm install \&\& npm run dev}
\end{enumerate}

If tables are missing, run \texttt{uv run alembic upgrade head} in \texttt{apps/api}.

\subsection{Notes}

\begin{itemize}[nosep]
  \item \texttt{.env}, \texttt{*.db} and local test-account files are git-ignored; each developer creates their own \texttt{.env}.
  \item On Windows use \texttt{127.0.0.1}, not \texttt{localhost}, for \texttt{AGENTFORGE\_API\_URL}.
  \item Keep \texttt{AUTH\_REQUIRED=false} for local use. SSO is currently paused.
  \item Tests: run \texttt{uv run pytest} in \texttt{apps/api}, \texttt{apps/core} and \texttt{apps/worker}.
\end{itemize}

% ------------------------------------------------------------------ 3
\section{Block types}

Every workflow needs one \emph{Start} and one \emph{End} block. Edges must form a
directed acyclic graph (no loops).

\begin{longtable}{@{}L{3.6cm}L{12cm}@{}}
\toprule
\textbf{Group} & \textbf{Blocks} \\
\midrule
\endhead
Triggers & Start, Manual, Webhook (own URL), Schedule (cron or interval), Chat \\
AI & \textbf{LLM} (one prompt, one answer); \textbf{Agent} (AI that uses tools over several steps, with memory, step limits and token budgets) \\
Logic & If, Switch, Merge, Wait, \textbf{Approval} (pauses until a person approves) \\
Data & Template (text with \texttt{\{\{ input.name \}\}}), Set (build fields), HTTP request, Company database \\
Integrations & Slack, GitHub, Notion, Airtable, Resend (email), Stripe (read-only), Google Sheets and Drive, Discord, HubSpot, Asana, Linear, Monday, ClickUp, Typeform, GitLab, Sentry, Vercel, Cloudflare (read-only), Webflow, Brandfetch, and search tools (Tavily, Serper, Exa, Firecrawl, Wikipedia) \\
End & Finishes the workflow and returns the result \\
\bottomrule
\end{longtable}

% ------------------------------------------------------------------ 4
\section{Pages of the application}

\begin{longtable}{@{}L{4.6cm}L{11cm}@{}}
\toprule
\textbf{Page} & \textbf{Purpose} \\
\midrule
\endhead
Home (\texttt{/}) & List of workflows, \textbf{+ New workflow}, and the role picker (Everyone, Sales, Marketing, Tech, Analyst, Finance) which tailors the examples shown in block help. \\
Editor (\texttt{/editor/<id>}) & Canvas to place and configure blocks, run, and use the Copilot. \\
Use (\texttt{/use/<id>}) & Simple page for non-technical users: inputs, one button, result as text or table. \\
Runs (\texttt{/graphs/<id>/runs}) & Run history with per-step details. \\
Triggers (\texttt{/graphs/<id>/triggers}) & Schedules and webhook addresses. \\
Integrations & Browse and connect the supported tools. \\
Credentials & Saved API keys and logins (Google uses OAuth). \\
Company data (\texttt{/data}) & Connect the company database read-only; browse tables; edit the data dictionary; try a query. \\
Files & Upload files that workflows can reference. \\
Settings & Team, Usage and quotas, Activity (audit log), Data retention. \\
\bottomrule
\end{longtable}

% ------------------------------------------------------------------ 5
\section{Getting started (about five minutes)}

\begin{enumerate}
  \item Sign in or sign up. With \texttt{AUTH\_REQUIRED=false} the app also works in single-user mode.
  \item On the Home page choose your \textbf{role} so that examples match your job.
  \item Press \textbf{+ New workflow}. A starter opens: Start $\rightarrow$ Template (``Hello \{\{ input.name \}\}!'') $\rightarrow$ End.
  \item Click a block to configure it. Each block has plain-language help with an example.
  \item Press \textbf{Run}, enter the input and read the result in the editor.
  \item Open \texttt{/use/<id>} to see what a non-technical colleague would see.
\end{enumerate}

% ------------------------------------------------------------------ 6
\section{Common tasks}

\subsection{Use AI in a workflow}
Save an API key under \emph{Credentials}. Add an \textbf{LLM} block for a single
answer, or an \textbf{Agent} block when the AI should use tools over several
steps. Agents cannot write to systems unless \texttt{allow\_writes} is enabled.

\subsection{Connect a tool (for example Slack)}
Open \emph{Integrations}, choose the tool, and add its key or sign in with OAuth.
Add the matching block to the workflow and select that credential.

\subsection{Connect the company database}
\begin{enumerate}[nosep]
  \item An owner opens \texttt{/data}. The page shows SQL to create a login that can only \emph{read}.
  \item The owner connects the database and presses \textbf{Test}; AgentForge lists the tables.
  \item Editors may browse tables and try a query. Nothing in the platform can change company data.
  \item Add a \textbf{Company database} block, or give an Agent the database tool.
\end{enumerate}

\subsection{Describe the data (data dictionary)}
In the dictionary editor, map business concepts (customer, invoice, loan, and so on)
to the company's real tables and columns, and add today's date, currency and
policies. The built-in \textbf{Treasury / Transaction Manager} role template is
written only in business terms, so it works on any company's schema; only the
dictionary differs.

\subsection{Run automatically}
Use a \textbf{Schedule} trigger (for example every Monday at 9) or a \textbf{Webhook}
trigger, which gives a URL that a CRM or form can call.

\subsection{Keep a human in control}
Add an \textbf{Approval} block. The run pauses and shows ``Waiting for someone to
approve'' until a person decides.

\subsection{Fix a failed run}
Open the run, find the failing step, and ask the Copilot to \emph{explain} or
\emph{fix} it.

\subsection{Versions and publishing}
Workflows keep version history with rollback, and can be duplicated, deleted and
restored. \textbf{Publish} makes a workflow available; \textbf{Activate} turns on
its triggers.

% ------------------------------------------------------------------ 7
\section{Roles and permissions}

\begin{longtable}{@{}L{2.6cm}L{13cm}@{}}
\toprule
\textbf{Role} & \textbf{Can do} \\
\midrule
\endhead
Owner & Everything, including connecting company databases, managing the team and transferring ownership. \\
Editor & Build and run workflows, browse company data, try queries. \\
Viewer & View workflows only; creation and edit controls are hidden. \\
\bottomrule
\end{longtable}

\subsection{Security and governance features}
\begin{itemize}[nosep]
  \item Credentials are stored encrypted with a master key (\texttt{CREDENTIALS\_MASTER\_KEY}).
  \item Company databases are accessed read-only; HTTP blocks accept only public destinations on ports 80 and 443 (SSRF protection).
  \item Audit log with verification and CSV export; usage quotas; data-retention settings.
  \item Rate limits on Copilot and data-source query endpoints.
  \item Instructions embedded in data (for example payment remarks) are treated as data, never as commands to an agent.
\end{itemize}

% ------------------------------------------------------------------ 8
\section{User journey}

\subsection{Persona A: Maya, Finance owner (builder)}

\begin{longtable}{@{}L{2.4cm}L{7.2cm}L{6cm}@{}}
\toprule
\textbf{Stage} & \textbf{What she does} & \textbf{Need / how the product helps} \\
\midrule
\endhead
1. Discover & Signs up, creates a workspace, picks the Finance role. & ``Is this for me?'' Role-based examples reassure her. \\
2. Set up & Adds an AI key under Credentials. Connects the company database on \texttt{/data} with a read-only login. Fills in the data dictionary. & Trust that data cannot be changed: read-only wording and the SQL snippet. \\
3. Build & Creates a workflow with an Agent block (Treasury template) and a Slack block, or asks the Copilot to draft it. & Speed: Copilot and per-block help. \\
4. Test & Presses Run and reads the result in the editor. & Visible steps and a clear result. \\
5. Automate & Adds a Schedule trigger and an Approval block before any sensitive action, then activates. & A human stays in control. \\
6. Govern & Reviews usage and quotas, the activity log, and invites teammates as editors or viewers. & Cost control and accountability. \\
\bottomrule
\end{longtable}

\subsection{Persona B: Raj, Operations user (consumer)}

\begin{longtable}{@{}L{2.4cm}L{13.2cm}@{}}
\toprule
\textbf{Stage} & \textbf{What he does} \\
\midrule
\endhead
1. Receive & Gets a link \texttt{/use/<id>} from Maya. \\
2. Fill in & Types a few boxes, such as customer name or month. \\
3. Run & Presses the button and sees progress (``Working\ldots{} step 2 of 3''). \\
4. Read & Gets a plain-text answer or a table; presses \textbf{Run again} to keep his inputs. \\
5. Escalate & If something looks wrong he tells Maya, who opens the run history and uses the Copilot's fix suggestion. \\
\bottomrule
\end{longtable}

\subsection{Persona C: Senior developer (operator)}

\begin{longtable}{@{}L{2.4cm}L{13.2cm}@{}}
\toprule
\textbf{Stage} & \textbf{What they do} \\
\midrule
\endhead
1. Install & Clones the repository and starts the stack with Docker (Section 2). \\
2. Verify & Opens the API docs at \texttt{:8000/docs} and the Temporal UI at \texttt{:8233}. \\
3. Explore & Creates a sample workflow, runs it, and inspects runs in the web app and in Temporal. \\
4. Review & Checks the credential vault, audit log, read-only data-source handling and test suites. \\
\bottomrule
\end{longtable}

\subsection{End-to-end flow}

\begin{center}
\small
Sign in $\rightarrow$ choose role $\rightarrow$ connect credentials / company data $\rightarrow$ build workflow (or ask Copilot)\\
$\rightarrow$ test run $\rightarrow$ publish and activate $\rightarrow$ colleagues use \texttt{/use/<id>} $\rightarrow$ monitor runs, usage and audit log
\end{center}

% ------------------------------------------------------------------ 9
\section{Troubleshooting}

\begin{longtable}{@{}L{5.6cm}L{10cm}@{}}
\toprule
\textbf{Symptom} & \textbf{Check} \\
\midrule
\endhead
Home shows ``API: not reachable'' & The API is not running; check port 8000 and \texttt{.env}. \\
Runs stay ``queued'' & The worker or Temporal is not running. \\
Copilot or AI blocks do not respond & AI key missing or invalid, or the rate limit was reached. \\
Database ``Test'' fails & Wrong host or password, or the read-only user lacks \texttt{SELECT}. \\
Run stuck ``waiting'' & An Approval block is waiting for a person. \\
Missing tables on first start & Run \texttt{uv run alembic upgrade head} in \texttt{apps/api}. \\
\bottomrule
\end{longtable}

\vfill
\begin{center}
\footnotesize\textit{This guide was derived from the source code (pages, block types, role templates and API routes).
Button labels and minor UI details may differ slightly from the running application.}
\end{center}

\end{document}
